% Paper A: "Localizing the Listen Layer in Speech LLMs"
% Status: LaTeX Skeleton v1.0 — all sections scaffolded, content blocks marked
% Created: 2026-03-07 (cycle Q058)
% Track: T3 (Listen vs Guess)
%
% HOW TO USE:
%   - Content blocks marked %[CONTENT] are ready to paste from paper-a-pitch.md
%   - %[RESULT] blocks are placeholders until experiments run
%   - %[FIG] blocks = figure placeholders
%   - Compile with: pdflatex main.tex (or overleaf)
%
% TARGET VENUE: NeurIPS 2026 (primary) / Interspeech 2026 (fallback)
% FORMAT: NeurIPS 2026 template (swap \documentclass if targeting Interspeech)

\documentclass{article}

% ─── Packages ──────────────────────────────────────────────────────────────────
\usepackage[final]{neurips_2026}   % TODO: swap to interspeech_2026 if needed
\usepackage[utf8]{inputenc}
\usepackage[T1]{fontenc}
\usepackage{hyperref}
\usepackage{url}
\usepackage{booktabs}
\usepackage{amsfonts}
\usepackage{amsmath}
\usepackage{amssymb}
\usepackage{nicefrac}
\usepackage{microtype}
\usepackage{graphicx}
\usepackage{subfigure}
\usepackage{multirow}
\usepackage{xcolor}
\usepackage{algorithm}
\usepackage{algorithmic}
\usepackage{todonotes}  % remove before final submission

% ─── Convenience macros ────────────────────────────────────────────────────────
\newcommand{\gc}{\text{gc}}
\newcommand{\iia}{\text{IIA}}
\newcommand{\Lstar}{L^*}
\newcommand{\etal}{\textit{et al.}}
\newcommand{\eg}{\textit{e.g.},}
\newcommand{\ie}{\textit{i.e.},}
\newcommand{\todo}[1]{\textcolor{red}{[TODO: #1]}}
\newcommand{\result}[1]{\textcolor{blue}{[RESULT: #1]}}
\newcommand{\figref}[1]{Figure~\ref{#1}}
\newcommand{\tabref}[1]{Table~\ref{#1}}
\newcommand{\secref}[1]{Section~\ref{#1}}

% ─── Title & Authors ───────────────────────────────────────────────────────────
\title{Localizing the Listen Layer in Speech Language Models}

\author{
  Leonardo Hu\thanks{Equal contribution.} \\
  Graduate Institute of Communication Engineering \\
  National Taiwan University \\
  \texttt{leo@ntu.edu.tw} \\
  \And
  \todo{Co-author: 智凱哥 (AudioLens, ASRU 2025) — confirm collaboration} \\
  \todo{Affiliation} \\
  \texttt{\todo{email}} \\
}

\begin{document}

\maketitle

% ─────────────────────────────────────────────────────────────────────────────
\begin{abstract}
% [CONTENT] Abstract v0.6 — paste from paper-a-pitch.md §Abstract
%
% Large audio-language models (LALMs) can answer questions about audio content,
% but it remains unclear *where* in their forward pass audio information becomes
% causally decisive. ...
%
% STATUS: v0.6 draft ready (paper-a-pitch.md §Abstract). Paste here.
\todo{Paste abstract v0.6 from paper-a-pitch.md. ~150 words. Safe to share.}
\end{abstract}

% ─────────────────────────────────────────────────────────────────────────────
\section{Introduction}
\label{sec:intro}

% [CONTENT] §1 Introduction — 3 paragraphs, LaTeX-ready in paper-a-pitch.md §v1.2
% Para 1: Problem motivation (behavioral dominance known, causal localization missing)
% Para 2: Prior work gap (SCD — Store-Contribute Dissociation; probing ≠ causal)
% Para 3: Contribution (Listen Layer, gc(L), DAS-IIT, experiments on Whisper + Qwen2-Audio)
%
% KEY CITES: ALME (Li et al. 2602.11488), Cascade Equivalence (2602.17598),
%   MiSTER-E, DashengTokenizer (2602.23765), AG-REPA (2603.01006),
%   Braun et al. 2025, Hase et al. 2023, AudioLens, FCCT (2511.05923),
%   Geiger et al. (2301.04709, 2303.02536), Sutter et al. (2507.08802),
%   Choi et al. (2602.18899), ALME (2602.11488), SpeechTokenizer (2506.04492),
%   Joshi et al. (2602.16698), MPAR² (2603.02266)

\todo{Paste §1 Introduction (3 paragraphs) from paper-a-pitch.md §v1.2/v1.9}

% ─────────────────────────────────────────────────────────────────────────────
\section{Related Work}
\label{sec:related}

% [CONTENT] §2 Related Work — 3 subsections, LaTeX-ready in paper-a-pitch.md §v1.3
%
% STRUCTURE:
%   2.1 Modality Grounding in Audio-Language Models
%   2.2 Mechanistic Interpretability of Audio and Multimodal Models
%   2.3 Causal Abstraction and Distributed Alignment Search
%
% TABLE 1: 5-paper comparison table (EmbedLens, Liu et al. 2025, FCCT, AudioLens, Paper A)
%   Columns: Paper | Modality | Method | Causal? | Grounded Metric?
%   Rows: see paper-a-pitch.md §Connections to Related Work

\todo{Paste §2 Related Work (3 subsections + Table 1) from paper-a-pitch.md §v1.3}

\begin{table}[h]
  \caption{Comparison of related work on audio/visual information localization.}
  \label{tab:related}
  \centering
  \begin{tabular}{lllcc}
    \toprule
    \textbf{Paper} & \textbf{Modality} & \textbf{Method} & \textbf{Causal?} & \textbf{Grounded Metric?} \\
    \midrule
    EmbedLens \cite{fan2026embedlens}       & Vision  & Probing (sink/dead/alive) & \xmark & \xmark \\
    Liu \etal\ 2025 \cite{liu2025kvflow}    & Vision  & KV-token flow (observational) & \xmark & \xmark \\
    FCCT \cite{li2026fcct}                   & Vision  & Causal tracing (vanilla) & \cmark & \xmark \\
    AudioLens \cite{audiolens2025}           & Speech  & Logit lens (observational) & \xmark & \xmark \\
    \textbf{Paper A (ours)}                  & Speech  & DAS-IIT \gc(L)            & \cmark\cmark & \cmark \\
    \bottomrule
  \end{tabular}
\end{table}

% ─────────────────────────────────────────────────────────────────────────────
\section{Method}
\label{sec:method}

% [CONTENT] §3 Method — 8 subsections, LaTeX-ready in paper-a-pitch.md §v1.4 + §v2.1
%
% STRUCTURE:
%   3.1 Task Formulation       (causal abstraction, gc(L) definition, L* definition)
%   3.2 Stimuli                (phonological minimal pairs + ALME conflict pairs + RVQ corruption)
%   3.3 Distributed Alignment Search (DAS algorithm, L_IIT loss, Cayley parametrization)
%   3.4 Direction Extraction   (diff-of-means MMProbe, PROBE_LAYER vs INTERVENE_LAYER sweep)
%   3.5 Decomposability Ablation (voicing ⊥ phoneme-identity subspace test)
%   3.6 Connector Subspace Transfer Test (R_encoder applied at LLM layer 0)
%   3.7 Experimental Setup     (Whisper-small + Qwen2-Audio-7B, NNsight, pyvene)
%   3.8 Evaluation Protocol    (bootstrap 95% CI, peak condition, baselines B1–B4)

\subsection{Task Formulation}
\label{sec:task}

% [CONTENT] §3.1 — paste from paper-a-pitch.md §v1.4 §3.1
% Key equations to typeset:
%   gc(L) = IIA_DAS(L; A)
%   L* = argmax_L gc(L)

\todo{Paste §3.1 Task Formulation}

\begin{equation}
  \gc(L) = \iia_{\text{DAS}}(L;\, A)
  \label{eq:gc}
\end{equation}

\begin{equation}
  \Lstar = \underset{L}{\arg\max}\; \gc(L)
  \label{eq:lstar}
\end{equation}

\subsection{Stimuli}
\label{sec:stimuli}

% [CONTENT] §3.2 — paste from paper-a-pitch.md §v1.4 §3.2
% Phase 1: Choi et al. phonological minimal pairs (voicing contrasts)
% Phase 2: ALME 57K conflict stimuli + RVQ-layer-selective corruptions (SpeechTokenizer)

\todo{Paste §3.2 Stimuli}

\subsection{Distributed Alignment Search}
\label{sec:das}

% [CONTENT] §3.3 — paste from paper-a-pitch.md §v1.4 §3.3
% Key equation to typeset (IIT loss):

\todo{Paste §3.3 DAS algorithm}

\begin{equation}
  \mathcal{L}_{\text{IIT}}(R) =
  \mathbb{E}_{(x^c, x^n, y^*)} \left[
    \ell\!\left(
      \mathcal{M}\!\left(x^n \,\big|\, h_L \leftarrow R^{-1} P R h_L^c \right),\, y^*
    \right)
  \right]
  \label{eq:iit_loss}
\end{equation}

\subsection{Direction Extraction}
\label{sec:direction}

% [CONTENT] §3.4 — paste from paper-a-pitch.md §v1.4 §3.4
% Difference-of-means estimator (MMProbe):

\todo{Paste §3.4 Direction Extraction}

\begin{equation}
  \mathbf{d}_{\text{voicing}}^L =
  \mathbb{E}\!\left[\mathbf{h}_L \mid \text{voiced}\right]
  - \mathbb{E}\!\left[\mathbf{h}_L \mid \text{unvoiced}\right]
  \label{eq:mmprobe}
\end{equation}

\subsection{Decomposability Ablation}
\label{sec:decomp}

% [CONTENT] §3.5 — paste from paper-a-pitch.md §v1.4 §3.5
% Test: R_voicing ⊥ R_phoneme at L*?
% decomp(L*) = 1 - |cos∠(R_voicing, R_phoneme)|

\todo{Paste §3.5 Decomposability Ablation}

\subsection{Connector Subspace Transfer Test}
\label{sec:connector}

% [CONTENT] §3.6 — paste from paper-a-pitch.md §v1.4 §3.6
% IIA_transfer = IIA(M_LLM, R_encoder, L_LLM=0)
% Three interpretations: volume-preserving / rotation-preserving / geometry-destroying

\todo{Paste §3.6 Connector Subspace Transfer Test}

\subsection{Experimental Setup}
\label{sec:setup}

% [CONTENT] §3.7 — paste from paper-a-pitch.md §v1.4 §3.7
% Models: Whisper-small (244M, 6 encoder layers) + Qwen2-Audio-7B (32 LLM layers)
% Tools: NNsight (NOT circuit-tracer — cross-attention), pyvene RotatedSpaceIntervention

\todo{Paste §3.7 Experimental Setup}

\subsection{Evaluation Protocol and Metrics}
\label{sec:eval}

% [CONTENT] §3.8 — paste from paper-a-pitch.md §v2.1
% 3.8.1 Primary Metric: gc(L) definition + L* peak condition (bootstrap 95% CI)
% 3.8.2 Bootstrap Protocol (1000 resamples, non-parametric)
% 3.8.3 Baseline Comparisons (B1 random-init DAS, B2 vanilla patching, B3 MFA, B4 trivial flat-gc)
% 3.8.4 Per-Experiment Evaluation Conditions (E1 passing/failure, E2 passing/failure)
% 3.8.5 Per-Class Reporting (risk A6 mitigation)

\todo{Paste §3.8 Evaluation Protocol (from paper-a-pitch.md §v2.1)}

% ─────────────────────────────────────────────────────────────────────────────
\section{Experiments and Results}
\label{sec:experiments}

% [CONTENT] §4 Experiments — 7 subsections, LaTeX-ready in paper-a-pitch.md §v1.5 + §v2.0 + §v2.2

\subsection{E1: Whisper-small gc(L) Curve}
\label{sec:e1_gc}

% [CONTENT] §4.1 — paste from paper-a-pitch.md §v1.5 §4.1
% Expected: gc(L) peaks at L ≈ 3 (50% encoder depth), Triple Convergence
% Cites: AudioSAE (Aparin), "Beyond Transcription" (Glazer), whisper_hook_demo.py, SCD

\todo{Paste §4.1 E1 gc(L) curve predictions}

%[FIG] Figure 1: gc(L) curve for Whisper-small (6 encoder layers)
%      x-axis: encoder layer L (0–5)
%      y-axis: gc(L) = DAS-IIA
%      error bars: bootstrap 95% CI (1000 resamples)
%      baselines: B1 (random-init DAS), B4 (flat-gc = 0.5)
%      annotation: L* marker at peak
\begin{figure}[h]
  \centering
  \fbox{\parbox{0.75\textwidth}{\centering
    \textbf{[FIGURE 1 PLACEHOLDER]}\\[1ex]
    \textit{gc(L) curve for Whisper-small encoder (6 layers).}\\
    x-axis: encoder layer $L$; y-axis: $\gc(L)$;\\
    error bars: bootstrap 95\% CI; baselines B1 (random-init) and B4 (flat).\\
    Expected: sharp peak at $L^* \approx 3$ ($\approx$50\% depth).
  }}
  \caption{Grounding coefficient \gc(L) across Whisper-small encoder layers (E1). 
    The Listen Layer $\Lstar$ is identified by the peak satisfying the bootstrap 
    CI isolation criterion (\S\ref{sec:eval}).}
  \label{fig:e1_gc_curve}
\end{figure}

\subsection{E1: Decomposability Ablation at L*}
\label{sec:e1_decomp}

% [CONTENT] §4.2 — paste from paper-a-pitch.md §v1.5 §4.2
% Expected: decomp(L*) ≈ 0.8–0.9 (voicing ⊥ phoneme-identity = abstract phonological encoding)

\todo{Paste §4.2 Decomposability ablation predictions}

%[FIG] Figure 2: 2D probe_layer × intervene_layer IIA heatmap (A4 check)
%      "lower-triangular stripe" near L* as predicted in paper-a-pitch.md §Q16
\begin{figure}[h]
  \centering
  \fbox{\parbox{0.75\textwidth}{\centering
    \textbf{[FIGURE 2 PLACEHOLDER]}\\[1ex]
    \textit{2D heatmap: probe\_layer $\times$ intervene\_layer IIA.}\\
    Expected: lower-triangular stripe near $L^*$ (\S\ref{sec:direction}).\\
    (Prediction: high IIA when probe\_layer $\leq L^*$ AND intervene\_layer $= L^*$.)
  }}
  \caption{Probe-layer × intervene-layer IIA heatmap for Whisper-small (E1).
    The lower-triangular stripe pattern validates that causal direction is extracted
    before $\Lstar$ and applied at $\Lstar$ (A4 assumption check, \S\ref{sec:related}).}
  \label{fig:e1_heatmap}
\end{figure}

\subsection{E1: Connector Subspace Transfer}
\label{sec:e1_connector}

% [CONTENT] §4.3 — paste from paper-a-pitch.md §v1.5 §4.3
% IIA_transfer outcomes: volume-preserving / rotation-preserving / geometry-destroying

\todo{Paste §4.3 Connector subspace transfer predictions}

%[TABLE] Table 2: Connector transfer results
%   Rows: (a) IIA_transfer (no retraining), (b) Retrained IIA at LLM layer 0
%   Columns: gc at encoder L*, IIA_transfer, Retrained IIA (L=0), Interpretation
\begin{table}[h]
  \caption{Connector subspace transfer test results (E1, Gap \#18).}
  \label{tab:connector}
  \centering
  \begin{tabular}{lccc}
    \toprule
    \textbf{Condition} & $\gc(L^*_\text{enc})$ & $\iia_\text{transfer}$ & Retrained IIA $(L=0)$ \\
    \midrule
    E1 Whisper-small & \result{---} & \result{---} & \result{---} \\
    \bottomrule
  \end{tabular}
\end{table}

\subsection{E2: Qwen2-Audio-7B gc(L) Sweep}
\label{sec:e2}

% [CONTENT] §4.4 — paste from paper-a-pitch.md §v1.5 §4.4
% Expected: gc(L) peaks at L ∈ {14–22} of LLM layers
% Cites: Zhao et al. ESN clustering (2601.03115), MPAR² (2603.02266), Heimersheim & Nanda 2024

\todo{Paste §4.4 E2 Qwen2-Audio-7B gc(L) prediction}

%[FIG] Figure 3 (was Figure 4 in pitch): Full gc(L) curve for Qwen2-Audio-7B
%      x-axis: LLM layer L (0–31) + encoder layers
%      y-axis: gc(L) = DAS-IIA
%      show encoder → connector → LLM boundary markers
\begin{figure}[h]
  \centering
  \fbox{\parbox{0.75\textwidth}{\centering
    \textbf{[FIGURE 3 PLACEHOLDER]}\\[1ex]
    \textit{gc(L) curve for Qwen2-Audio-7B (32 LLM layers + Whisper encoder).}\\
    x-axis: layer index; boundaries: Whisper encoder | connector | LLM backbone.\\
    Expected: peak at $L \in \{14$–$22\}$ (middle-dominant SCD pattern).
  }}
  \caption{Grounding coefficient \gc(L) across Qwen2-Audio-7B (E2).
    The \textit{middle-dominant} profile (peak at mid-LLM depth, drop at upper layers)
    indicates Tier 3 grounding (text-prior dominance at late layers).}
  \label{fig:e2_qwen}
\end{figure}

\subsection{Predicted Failures and Contingencies}
\label{sec:contingencies}

% [CONTENT] §4.5 — paste from paper-a-pitch.md §v1.5 §4.5
% Risks A1–A6 with contingency plans

\todo{Paste §4.5 Predicted failures and contingencies}

\subsection{A Taxonomy of Grounding Failures}
\label{sec:taxonomy}

% [CONTENT] §4.6 — paste from paper-a-pitch.md §v2.0
% 3-Tier Model: Tier 1 (Codec), Tier 2 (Connector Bottleneck), Tier 3 (LLM Modality Collapse)
% Each tier: mechanism + empirical support + gc(k) signature + diagnostic + implication

\todo{Paste §4.6 Grounding Failure Taxonomy (3-Tier Model) from paper-a-pitch.md §v2.0}

%[TABLE] Table 3: Falsifiability summary — per-tier gc(k) signatures
%   Columns: Tier | Failure Site | gc(k) Pattern | enc gc? | LLM L0? | L_mid? | L_late?
\begin{table}[h]
  \caption{Grounding failure taxonomy: gc(L) diagnostic signatures by tier.}
  \label{tab:taxonomy}
  \centering
  \small
  \begin{tabular}{llccccc}
    \toprule
    \textbf{Tier} & \textbf{Failure Site} & \textbf{gc(L) Pattern}
      & \textbf{Enc $\gc(\Lstar)$}
      & \textbf{LLM $L_0$}
      & \textbf{$L_\text{mid}$}
      & \textbf{$L_\text{late}$} \\
    \midrule
    \textbf{T1} & Codec (RVQ)        & Flat ($\approx$chance)       & $\approx$ch. & $\approx$ch. & $\approx$ch. & $\approx$ch. \\
    \textbf{T2} & Connector          & Cliff at boundary             & HIGH         & $\approx$ch. & $\approx$ch. & $\approx$ch. \\
    \textbf{T3} & LLM (text prior)   & Mid-peak + late drop          & HIGH         & Med.         & HIGH         & DROP         \\
    \textbf{None} & No failure       & Rising plateau                & HIGH         & Med.         & HIGH         & HIGH         \\
    \bottomrule
  \end{tabular}
\end{table}

\subsection{Grounding Failure Diagnostic Protocol}
\label{sec:diagnostic}

% [CONTENT] §4.7 — paste from paper-a-pitch.md §v2.2
% Sequential protocol: Step 1 (Codec Probe) → Step 2 (Encoder gc sweep) →
%   Step 3 (Connector Transfer) → Step 4 (LLM gc sweep)
% Table 4: per-tier diagnostic tests + thresholds + time cost

\todo{Paste §4.7 Diagnostic Protocol from paper-a-pitch.md §v2.2}

%[TABLE] Table 4: Per-tier diagnostic tests
%   Columns: Tier | Failure Site | Diagnostic Test | Passing Threshold | Time Cost
\begin{table}[h]
  \caption{Diagnostic protocol: per-tier tests, thresholds, and time cost.}
  \label{tab:diagnostic}
  \centering
  \small
  \begin{tabular}{llp{4.5cm}p{2.5cm}r}
    \toprule
    \textbf{Tier} & \textbf{Site} & \textbf{Diagnostic Test}
      & \textbf{Threshold} & \textbf{Cost} \\
    \midrule
    T1a & Codec (RVQ)           & SpeechTokenizer L1 reconstruction preserves contrast? & probe acc. $<$ 0.55 & 2 min, CPU \\
    T1b & Codec (encoder gc)    & max encoder $\gc(L)$ & $<$ 0.55 everywhere & 30 min, CPU \\
    T2a & Connector (transfer)  & $\iia_\text{transfer}$ at LLM layer 0 & $<$ 0.55 & 5 min, CPU \\
    T2b & Connector (retrained) & DAS retrained at LLM layer 0 & retrained IIA $<$ 0.55 & 10 min, CPU \\
    T3  & LLM late-layer        & $\gc(L_\text{late})$ drop vs $\gc(L_\text{mid})$ & drop $\geq$ 0.10 (CI non-overlapping) & 1 day, GPU \\
    None & No failure           & $\gc(L_\text{late}) \geq \gc(L_\text{mid}) - 0.05$ & plateau holds & 1 day, GPU \\
    \bottomrule
  \end{tabular}
\end{table}

% ─────────────────────────────────────────────────────────────────────────────
\section{Discussion}
\label{sec:discussion}

% [CONTENT] §5 skeleton — paste from paper-a-pitch.md §v1.6
% PROSE BLOCKED until experimental results. Use stubs only.
%
% STRUCTURE:
%   5.1 The Listen Layer as a Causal Bottleneck (SCD confirmation)
%   5.2 Phonological Abstraction vs. Table-Lookup at L*
%   5.3 What the Connector Does to Phonological Geometry (Gap #18)
%   5.4 Grounding Profile Across Generation vs. Understanding
%   5.5 Limitations and Future Work

\subsection{The Listen Layer as a Causal Bottleneck}
\label{sec:disc_bottleneck}

% [CONTENT TRIGGER: Figure 1 gc(L) curve + Figure 2 heatmap — WRITE AFTER E1]
\todo{Write §5.1 after E1 results. Key: SCD confirmation (storage ≠ causal contribution).
  If decomp(L*) confirms abstract phonological encoding: discuss Store-Contribute Dissociation.}

\subsection{Phonological Abstraction vs.\ Table-Lookup at $\Lstar$}
\label{sec:disc_abstraction}

% [CONTENT TRIGGER: Table 1 phono-init vs random-init DAS ablation — WRITE AFTER E1]
\todo{Write §5.2 after E1 results. Key: if decomp(L*) ≈ 0.8–0.9 → abstract encoding confirmed;
  contrast with text LLM factual recall (Meng et al. 2022).}

\subsection{What the Connector Does to Phonological Geometry}
\label{sec:disc_connector}

% [CONTENT TRIGGER: Table 2 connector subspace transfer results — WRITE AFTER E1]
\todo{Write §5.3 after Gap \#18 experiment. Key: connector as information bottleneck
  vs. geometry-preserving rotation. Implications for Modality Collapse (2602.23136).}

\subsection{Grounding Profile Across Generation vs.\ Understanding}
\label{sec:disc_gen_vs_under}

% [CONTENT TRIGGER: Figure 3 Qwen2-Audio gc(L) sweep — WRITE AFTER E2]
\todo{Write §5.4 after E2. Key: middle-dominant (understanding) vs. early-dominant (generation,
  AG-REPA 2603.01006) SCD asymmetry. Implications for interpretability method choice.}

\subsection{Limitations and Future Work}
\label{sec:limitations}

% [CONTENT] §5.5 stub from paper-a-pitch.md §v1.6 — expand before final submission
\todo{Paste §5.5 stub from pitch. Key limitations:
  (1) linear alignment assumption (DAS) — non-linear connectors may need kernel DAS (Gap \#25);
  (2) phonological features tested = voicing/manner/place only;
  (3) Whisper-small + Qwen2-Audio-7B only.
  Future: AudioSAEBench (Paper B), LoRA-adapted models (Track 4), adversarial audio (Track 5).}

% ─────────────────────────────────────────────────────────────────────────────
\section{Conclusion}
\label{sec:conclusion}

\todo{Write conclusion after §5 discussion prose is complete (post-E1/E2).
  Key: (1) introduced Listen Layer + gc(L); (2) causal DAS-IIT localization;
  (3) 3-tier grounding failure taxonomy; (4) diagnostic protocol.}

% ─────────────────────────────────────────────────────────────────────────────
% ─── Bibliography ─────────────────────────────────────────────────────────────
\bibliographystyle{plainnat}
\bibliography{refs}

% ─────────────────────────────────────────────────────────────────────────────
% ─── Appendix ─────────────────────────────────────────────────────────────────
\appendix

\section{Supplementary Figures}
\label{app:figures}

%[FIG App-A] Per-phoneme-class gc(L) curves (risk A6 mitigation)
\begin{figure}[h]
  \centering
  \fbox{\parbox{0.75\textwidth}{\centering
    \textbf{[APPENDIX FIGURE A PLACEHOLDER]}\\[1ex]
    \textit{Per-phoneme-class gc(L) curves.}\\
    Classes: [b]/[p] voicing, [d]/[t] voicing (+ rare contrasts if available).\\
    Expected: common phonemes maintain $\gc(L_\text{late})$; rare phonemes drop.
  }}
  \caption{Per-phoneme-class grounding coefficient curves (§3.8.5 — risk A6 mitigation).
    Divergence between common and rare phoneme classes would indicate class-conditional
    grounding failure at upper LLM layers.}
  \label{fig:app_per_class}
\end{figure}

\section{Statistical Details}
\label{app:stats}

% Bootstrap protocol details, seed, number of resamples

\todo{Add bootstrap protocol details: 1000 resamples, seed=42,
  non-parametric CI, peak isolation criterion (§3.8.2).}

\section{DAS Implementation Details}
\label{app:das}

% pyvene RotatedSpaceIntervention config, Cayley parametrization, optimizer

\todo{Add pyvene implementation details, hyperparameters, training setup.}

\section{Dataset Details}
\label{app:data}

% Choi et al. minimal pairs: 96 languages, voicing contrasts
% ALME: 57K stimuli, split details
% SpeechTokenizer RVQ-selective corruption construction

\todo{Add dataset details, splits, preprocessing.}

\end{document}
